\subsection{A rate proportional to the current amount}

A simple differential model begins with a law of change rather than a ready-made formula.

\begin{corebox}[Proportional growth law]
If a quantity $P(t)$ changes at a rate proportional to its current value, then
\[
P'(t)=kP(t).
\]
Positive $k$ models growth and negative $k$ models decay.
\end{corebox}

\begin{examplebox}[Checking an exponential model]
Let
\[
P(t)=P_0e^{kt}.
\]
Then
\[
P'(t)=kP_0e^{kt}=kP(t),
\]
so the exponential function satisfies the proportional-rate law.
\end{examplebox}

\begin{interpretationbox}[The derivative states the mechanism]
The equation $P'=kP$ says that the present rate depends on the present amount. The exponential formula is a function whose derivative reproduces precisely that mechanism.
\end{interpretationbox}

The aim here is not a general course in differential equations. The example shows how a derivative can encode a modelling assumption and how a proposed function can be tested against it.
